Write \(\mathcal O_a=\{O\in\mathcal O:O\subseteq\mathcal U_a\}\) and \(m_a=|\mathcal O_a|\). Values of \(s_a\) outside \(\mathcal O_a\) are irrelevant. If \(k_a\le M\), define a total preorder by \(O\preceq_a O'\) if and only if \(s_a(O)\le s_a(O')\). Then
\[
s_a(O_{T,a})\le Q_{k_a}(s_a(O_{1,a}),\ldots,s_a(O_{M,a}))
\quad\Longleftrightarrow\quad
\#\{i:s_a(O_{i,a})<s_a(O_{T,a})\}<k_a.
\tag{1}
\]
Thus the joint rank event depends only on the two induced total preorders, including their tie classes. Conversely, if a preorder has ordered tie classes \(C_{a,1}\prec_a\cdots\prec_a C_{a,d_a}\), assign score zero when \(d_a=1\), and otherwise assign \((j-1)K/(d_a-1)\) to each orbit in \(C_{a,j}\). This realizes the preorder in \([0,K]\) using \(d_a\le m_a\) levels. If \(k_a=M+1\), the armwise event is identically true by the convention \(Q_{k_a}=+\infty\), so the same finite reduction remains valid. Since there are only finitely many total preorders, the maximum in \(\texttt{def:coverage-frontier}\) is attained by enumerating all preorder pairs. This proves exactness over the definition's stated class of bounded orbit-constant maps; it makes no enlargement to all invariant latent laws.

For evaluability, define the proof intermediates
\[
\kappa_{w,a}(O\mid z)=\sum_{v:(v,a)\in O}\kappa_{w,a}(v\mid z)
\]
and
\[
R_w(O_0,O_1)=\sum_z\nu(z)\kappa_{w,0}(O_0\mid z)\kappa_{w,1}(O_1\mid z).
\tag{2}
\]
The auxiliary-product assumption makes the two selector draws conditionally independent given their shared assignment, so (2) is their joint orbit law. The orbit pairs are independent and identically distributed across calibration replicates. If \(O_a(v)\) denotes the orbit containing \((v,a)\), then the joint rank-event probability for a fixed preorder pair equals the finite sum
\[
\begin{aligned}
&\sum_{v\in\mathcal V}q_{\mathrm{tar}}(v)
\sum_{((O_{i,0},O_{i,1}))_{i=1}^M\in(\mathcal O_0\times\mathcal O_1)^M}
\left\{\prod_{i=1}^M R_w(O_{i,0},O_{i,1})\right\}\\
&\qquad\times\prod_{a=0}^1\mathbf 1\!\left\{s_a(O_a(v))\le Q_{k_a}(s_a(O_{1,a}),\ldots,s_a(O_{M,a}))\right\}.
\end{aligned}
\tag{3}
\]
This expression preserves both the common target member and the within-replicate cross-arm coupling; no target-arm or calibration-arm independence has been introduced.

Let \(n=m_0m_1\). The \(n^M\) configuration weights in (3) may first be tabulated. The \(n\)-ary strings of length \(M\) admit a reflected enumeration in which successive strings differ in one coordinate: recursively enumerate the last \(M-1\) coordinates forward for even values of the first coordinate and backward for odd values. This construction proves the property by induction on \(M\). For a fixed target and preorder pair, initialize the two strict-lower counts in (1) in \(O(M)\) operations and update them in constant time after each one-coordinate change. When \(n\ge2\), \(M\le2^M\le n^M\); when \(n=1\), the sum is immediate. Hence (3), over all targets and preorder pairs, is evaluated in
\[
O\!\left(F_{m_0}F_{m_1}|\mathcal V|(m_0m_1)^M\right)
\tag{4}
\]
arithmetic operations after the joint law is tabulated, where \(F_m\) is the number of total preorders on an \(m\)-element labelled set.

It remains to prove that binary maps do not always attain the maximum. Consider four pair blocks with group \(S_2^4\) and the invariant assignment law that treats exactly one member per pair, equiprobably within each pair. Let the selector library contain one rule which, in either arm, chooses a block with probabilities
\[
q=\left(\frac{13}{20},\frac1{20},\frac14,\frac1{20}\right)
\]
and selects that block's unique observed member in the requested arm. The rule is observable and equivariant. Its two armwise auxiliary choices are conditionally independent and their block probabilities do not depend on the realized within-pair labels, so \(R_w(b_0,b_1)=q_{b_0}q_{b_1}\). Let
\[
q_{\mathrm{tar}}(v_{b,r})=\frac{p_b}{2},\qquad p=\left(\frac1{20},\frac1{10},\frac14,\frac35\right).
\]
Take \(M=3\), \(\alpha=1/2\), and \(\varepsilon_0=\varepsilon_1=1/4\), so \(k_0=k_1=3\). Put \(s_0\equiv0\), and order the arm \(1\) blocks as \(1\prec3\prec2\prec4\), represented by \(0,K/3,2K/3,K\). Arm one fails precisely when all three calibration blocks are strictly below the target block. Therefore its failure probability is
\[
F_{\mathrm{multi}}=\frac14\left(\frac{13}{20}\right)^3+\frac1{10}\left(\frac{18}{20}\right)^3+\frac35\left(\frac{19}{20}\right)^3=\frac{104957}{160000}.
\tag{5}
\]
Arm zero never fails, so the frontier objective equals
\[
F_{\mathrm{multi}}-\alpha=\frac{24957}{160000}.
\tag{6}
\]

For a binary map in arm \(a\), let \(H_a\subseteq\{1,2,3,4\}\) be its high-score blocks and put \(A_a=(1-q(H_a))^3\). Conditional on the common target block, inclusion--exclusion and the product form of \(R_w\) give the probability that at least one arm fails:
\[
F_{\mathrm{bin}}(H_0,H_1)=p(H_0)A_0+p(H_1)A_1-p(H_0\cap H_1)A_0A_1.
\tag{7}
\]
In the row order
\[
\varnothing,1,2,3,4,12,13,14,23,24,34,123,124,134,234,1234,
\]
direct rational evaluation of the sixteen choices for the other arm gives the following row maxima, all over denominator \(320000000\):
\[
\begin{split}
(&164616000,165302000,192052000,198366000,207906936,165912000,164712000,165787368,\\
&203032000,208582976,208582976,164632000,165793875,164631423,202894331,164616000).
\end{split}
\tag{8}
\]
Thus (8) is a finite certificate exhausting all \(16^2\) binary pairs, and
\[
\max_{H_0,H_1}F_{\mathrm{bin}}(H_0,H_1)=\frac{3259109}{5000000}.
\]
After subtracting \(\alpha=1/2\), the largest binary objective is \(759109/5000000\). Comparing it with (6) gives the strict gap
\[
\frac{24957}{160000}-\frac{759109}{5000000}=\frac{83189}{20000000}>0.
\tag{9}
\]
This proves nonattainment by binary score pairs in an admissible symmetric member of the asymmetric family.

Finally, suppose target and calibration orbit supports are disjoint in arm \(a\) and \(k_a\le M\). Assign score \(K\) to target-supported orbits of that arm and zero to every other orbit, and assign score zero in the other arm. The arm \(a\) target score is always \(K\), while all its calibration scores are zero, so its rank event is impossible. The joint event in (3) has probability zero, and the frontier objective is \(1-\alpha\). Since no objective can exceed \(1-\alpha\),
\[
\Gamma_{M,\alpha,\boldsymbol\varepsilon}(w)=1-\alpha.
\]
If \(k_a=M+1\), then \(Q_{k_a}=+\infty\), so that arm's indicator in (3) is identically one. Maximization therefore reduces exactly to the other arm's rank event. If both ranks equal \(M+1\), the joint event has probability one for every score pair and the frontier is zero. On the diagonal, the identities in \(\texttt{def:coverage-frontier}\), \(\texttt{def:finite-rank-error}\), and \(\texttt{def:rank-inversion}\) recover the symmetric frontier statement.